\title{FlashAttention-3: Fast and Accurate Attention with Asynchrony and Low-precision}

\begin{abstract}
The attention mechanism is a fundamental component within the widespread Transformer framework and serves as a primary limitation in scaling large language models and enabling extended-context scenarios. \textsc{FlashAttention} previously introduced a method to accelerate attention operations on GPUs by reducing redundant memory access. Despite these advancements, newer hardware features remain underutilized; for instance, \textsc{FlashAttention-2} utilizes only 35\% of the H100 GPU’s capabilities. To address this, we propose three strategies specifically tailored for Hopper GPUs: harnessing the asynchronous execution potential of Tensor Cores and TMA by (1) employing warp-specialization to overlap computations with data transfers, (2) merging block-level matrix multiplication and softmax computation, and (3) implementing block quantization together with non-coherent processing that takes advantage of FP8 precision hardware capabilities. Our solution, \textsc{FlashAttention-3}, demonstrates 1.5-2.0$\times$ faster attention throughput on H100 GPUs; FP16 performance reaches up to 740 TFLOPs/s (equivalent to 75\% hardware utilization), while FP8 approaches nearly 1.2 PFLOPs/s. Additionally, our FP8 \textsc{FlashAttention-3} technique reduces numerical error by a factor of 2.6 relative to a standard FP8 attention baseline.
\end{abstract}

\section{Introduction}
\label{sec:intro}

Within the Transformer framework~\citep{vaswani2017attention}, the chief computational limitation arises from the attention mechanism, as the calculation of self-attention scores between queries and keys exhibits quadratic complexity with respect to sequence length. Enabling attention mechanisms to effectively handle extended contexts would greatly expand their potential, allowing for sophisticated modeling and inference across multiple lengthy documents~\citep{guo2021longt5,shaham2022scrolls,peng2023yarn}, comprehensive analysis of files in substantial code repositories~\citep{roziere2023code, li2023starcoder}, as well as support for emerging modalities including high-definition imagery~\citep{chen2022scaling}, audio~\citep{gulati2020conformer}, and video~\citep{ho2022video}. Such advancements would facilitate novel uses, for example, maintaining prolonged histories in user interactions~\citep{sun2019bert4rec} or enabling agents to operate over extended procedural horizons~\citep{yao2022react}. Consequently, there has been considerable focus on accelerating the attention computation, particularly for longer contexts—whether by leveraging approximation techniques~\citep{katharopoulos2020transformers,choromanski2020rethinking,tay2020efficient}, implementing software-level optimizations~\citep{rabe2021self,dao2022flashattention,kwon2023efficient}, or by exploring fundamentally different model architectures~\citep{peng2023rwkv,sun2023retentive,gu2023mamba}.

This study extends the prior efforts of \citet{dao2022flashattention}, who formulated exact-attention methods by thoughtfully considering the GPU’s execution paradigms and inherent hardware constraints in their algorithmic architecture. In their seminal work, \citet{dao2022flashattention} proposed \textsc{FlashAttention}, an innovative tiling methodology for attention that consolidates all computations within a single GPU kernel, thereby removing the need for slow intermediate transfers to global memory. Subsequently, \citet{dao2023flashattention2} redesigned the technique as \textsc{FlashAttention-2}, enabling parallelization along the sequence length axis and optimizing the forward pass by processing blocks of the key and value matrices. This modification enhanced GPU workload distribution and occupancy. Despite these advances, our analysis indicates that \textsc{FlashAttention-2} still demonstrates suboptimal efficiency on advanced GPUs compared to highly tuned matrix-multiplication (GEMM) operations—reporting utilization rates of 35\% versus 80–90\% on the Hopper H100 device. Some of this inefficiency stems from implementation specifics, notably the use of instructions designed for earlier GPU generations (Ampere) instead of Hopper-tailored operations, especially in contexts exploiting Tensor Cores. Additional research, including ThunkerKitten~\citep{spector2024thunder} and cuDNN 9~\citep{cudnn9}, confirms that leveraging Hopper-specific instructions along with tile-centric abstractions both accelerates attention mechanisms and makes their development more straightforward.

At a more basic level, the core algorithm in \textsc{FlashAttention-2} is structured around a straightforward, synchronous model, and it does not incorporate explicit mechanisms for asynchrony or leverage low-precision arithmetic in its approach. Asynchrony arises due to specialized hardware that expedites the principal operations in machine learning tasks: dedicated units such as Tensor Cores for matrix multiplication and the Tensor Memory Accelerator (TMA) for efficient memory access, distinct from the general CUDA cores responsible for logic operations and both integer and floating point arithmetic. The progression to reduced precision formats—including FP8 on Hopper, FP4 on Blackwell, following FP16 introduced on Pascal in 2017 and BF16 on Ampere in 2020—has demonstrably allowed for two to four times greater throughput at equivalent power usage and chip area. In \cref{subsec:hardware}, we examine these advancements as enabled by Hopper hardware. The main technical obstacle lies in adapting the design of \textsc{FlashAttention-2} to harness such hardware capabilities: integrating asynchrony mandates careful structuring to allow concurrent computation between matrix multiplication and softmax, even though there is a dependency between their outputs; employing low-precision techniques, on the other hand, necessitates careful management of quantization error—this is particularly important for mitigating issues arising from outlier features found in LLMs~\citep{dettmers2208llm, sun2024massive}.

With this objective in mind, we introduce \textsc{FlashAttention-3}, which integrates and advances three novel concepts aimed at enhancing performance on modern GPU architectures.\footnote{Our findings are discussed with respect to NVIDIA's Hopper architecture, though the proposed algorithm remains effective on any GPU platform that supports strong asynchronous execution and low-precision operations.}

\begin{enumerate}

\item \textbf{Producer-Consumer asynchrony:} We introduce a warp-specialized software pipelining strategy that takes advantage of asynchronous execution between data movement and Tensor Core computations. This is achieved by assigning separate warps to data producers and consumers, thereby enhancing the algorithm’s capability to mask both memory access and instruction issue delays.

\item \textbf{Hiding softmax under asynchronous block-wise GEMMs:} We interleave the relatively slower, non-GEMM tasks associated with softmax—such as floating point multiply-add and exponential calculations—with asynchronous WGMMA-based GEMM operations. In doing so, we modify the \textsc{FlashAttention-2} algorithm to bypass inherent sequential dependencies between softmax and GEMMs. For instance, in the two-stage design, softmax processes one block of the scores matrix while simultaneously, WGMMA is asynchronously employed to compute the following block.

\item \textbf{Hardware-accelerated low-precision GEMM:} We revise the forward pass approach to leverage FP8 Tensor Cores for GEMM computation, leading to an approximate twofold increase in sustained TFLOPs/s. This transition necessitates aligning the memory layout expectations of WGMMA regarding FP32 accumulator and FP8 operand blocks. Through block quantization and incoherent processing, we alleviate the precision loss associated with adopting FP8 representation.
\end{enumerate}

To assess the practical effectiveness of our approach, we evaluate \textsc{FlashAttention-3} on the H100 SXM5 GPU across diverse parameter settings and demonstrate the following: (1) FP16 attains a speedup of 1.5-2.0$\times$ compared to \textsc{FlashAttention-2} during the forward pass, achieving performance up to 740 TFLOPs/s, and a 1.5-1.75$\times$ improvement during the backward pass; (2) FP8 reaches nearly 1.2 PFLOPs/s; and (3) for long input sequences, FP16 delivers superior results and FP8 remains highly competitive\footnote{Specifically, for head dimension 64, FP8 \textsc{FlashAttention-3} surpasses competitors, while for head dimensions 128 and 256, its performance matches state-of-the-art for non-causal scenarios and lags slightly for causal masking.} relative to the leading attention implementation in NVIDIA's cuDNN library. Furthermore, we confirm that FP16 \textsc{FlashAttention-3} preserves the same numerical accuracy as \textsc{FlashAttention-2} and outperforms traditional attention implementations since it computes intermediate values (such as softmax rescaling) in FP32. In addition, FP8 \textsc{FlashAttention-3} utilizing block quantization and incoherent computation proves to be 2.6$\times$ more accurate than the typical attention mechanism employing per-tensor quantization, particularly in the presence of outlier features.

We have released \textsc{FlashAttention-3} under a permissive license\footnote{\textsc{FlashAttention-3} can be accessed at \texttt{https://github.com/Dao-AILab/flash-attention}}, and intend to incorporate it into PyTorch and Hugging Face libraries in order to maximize accessibility and utility for researchers and developers.

\section{Background: Multi-Head Attention and GPU Characteristics}
\label{sec:background}

\subsection{Multi-Head Attention}
\label{subsec:multi_head_attn}

Let $\mathbf{Q}, \mathbf{K}, \mathbf{V} \in \mathbb{R}^{N \times d}$ be the query, key and value input sequences associated to a single head, where $N$ is the sequence length and $d$ is the head dimension. Then the attention output $\mathbf{O}$ is computed as:

\begin{equation*}
  \mathbf{S} = \alpha \mathbf{Q} \mathbf{K}^\top \in \mathbb{R}^{N \times N}, \quad \mathbf{P} = \mathrm{softmax}(\mathbf{S}) \in \mathbb{R}^{N \times N}, \quad \mathbf{O} = \mathbf{P}\mathbf{V} \in \mathbb{R}^{N \times d},
\end{equation*}

Here, the $\mathrm{softmax}$ operation is performed across each row, and the scaling factor $\alpha$ is usually chosen as $1/\sqrt{d}$. To mitigate issues with numerical stability caused by the exponential function, it is common to subtract $\mathrm{rowmax}(\mathbf{S})$ from $\mathbf{S}$ prior to applying the softmax. In the case of multi-head attention (MHA), separate query, key, and value projection matrices are used for each head. The associated computations can be efficiently parallelized over various heads and batches, yielding the complete output tensor.

Now let $\phi$ be a scalar loss function and let $\mathbf{d}(-) = \partial \phi / \partial (-)$ be notation for the gradient. Given the output gradient $\mathbf{dO} \in \mathbb{R}^{N \times d}$, we compute $\mathbf{dQ}$, $\mathbf{dK}$, and $\mathbf{dV}$ according to the chain rule as follows:

\begin{align*}
  \mathbf{dV} &= \mathbf{P}^\top \mathbf{dO} \in \mathbb{R}^{N \times d} \\
  \mathbf{dP} &= \mathbf{dO} \mathbf{V}^\top \in \mathbb{R}^{N \times N} \\
  \mathbf{dS} &= \mathrm{dsoftmax} (\mathbf{dP}) \in \mathbb{R}^{N \times N} \\
  \mathbf{dQ} &= \alpha \mathbf{dS} \mathbf{K} \in \mathbb{R}^{N \times d} \\
  \mathbf{dK} &= \alpha \mathbf{dS}^\top \mathbf{Q} \in \mathbb{R}^{N \times d},
\end{align*}

In this context, $\mathbf{d}s = (\mathrm{diag}(p) - p p^\top)\mathbf{d}p$ holds when $p = \mathrm{softmax}(s)$ with respect to the vector $s$. We denote $\mathrm{dsoftmax}(\mathbf{dP})$ as applying this expression to each row individually. Additionally, this operation remains efficiently parallelizable over both the number of heads and batches during the backward propagation phase of MHA.

\subsection{GPU hardware characteristics and execution model}
\label{subsec:hardware}

We outline the components of the GPU execution model that pertain to \textsc{FlashAttention-3}, centering our discussion on the NVIDIA Hopper architecture as a specific example of this model in practice.

\paragraph{Memory hierarchy:} On GPUs, the organization of memory follows a hierarchical structure, where data locales with higher capacity typically exhibit lower bandwidth (\cref{tab:gpu-hierarchy})\footnote{\citet{luo2024benchmarking} provides a value of 128 bytes per clock cycle per SM for shared memory bandwidth; our calculation multiplies this by 132 SMs and a boost clock rate of 1830 MHz.}. The off-chip DRAM, often referred to as HBM or global memory (GMEM), is uniformly accessible across all streaming multiprocessors (SMs). GMEM contents are automatically stored in the on-chip L2 cache to facilitate faster data retrieval. Each SM, in turn, is equipped with its own relatively small, multi-banked, on-chip cache—shared memory (SMEM)—which is directly managed by the programmer. At the innermost level, every SM houses a dedicated register file for fastest access.

\paragraph{Thread hierarchy:} In the GPU programming paradigm, execution units known as threads are systematically grouped to form a hierarchical structure. Starting from the smallest unit, this hierarchy consists of individual threads, followed by warps—each containing 32 threads. Next are warpgroups, which aggregate 4 consecutive warps. Further up are threadblocks, also referred to as cooperative thread arrays (CTAs), threadblock clusters (specific to Hopper architecture), and ultimately, grids at the highest level.

The two hierarchies exhibit a strong interconnection. All threads belonging to a single CTA are executed concurrently on the same SM, while CTAs grouped within a cluster are assigned together to a single GPC. SMEM is accessible to every thread in a CTA, while individual threads have exclusive access to up to 256 private registers (RMEM).

\paragraph{Asynchrony and warp-specialization:}

GPUs function as throughput-oriented processors, utilizing parallelism and asynchronous operations to mask delays in memory access and execution. Hopper introduces the Tensor Memory Accelerator (TMA), a specialized hardware component, to facilitate asynchronous data transfer between GMEM and SMEM \cite[\S7.29]{cuda}. In addition, Hopper’s Tensor Core—accessible through the warpgroup-wide WGMMA instruction \cite[\S9.7.14]{ptx}—operates asynchronously and can obtain its input data straight from shared memory, a capability not present in earlier designs like Ampere.

The presence of hardware mechanisms for asynchrony enables the implementation of warp-specialized kernels, wherein individual warps within a CTA are assigned distinct responsibilities as either producers or consumers, exclusively handling either data transfer or arithmetic operations, respectively. This separation generally enhances the compiler’s capacity to arrange instructions more efficiently \citep{warp-specialization-2011}. Furthermore, Hopper introduces a feature for runtime register redistribution among warpgroups, using the \verb|setmaxnreg| directive \cite[\S9.7.17.1]{ptx}, which permits allocating a larger portion of RMEM to warps executing MMAs, as opposed to those merely handling TMA operations (which can be executed by a single thread).

\paragraph{Low-precision number formats:}
\label{sec:low-precision-gpu}
Contemporary GPUs are equipped with dedicated hardware components designed to expedite computations at reduced numerical precision. As an illustration, the WGMMA instruction enables access to FP8 Tensor Cores on Hopper, resulting in twice the throughput per SM in comparison to computations performed with FP16 or BF16 formats.

To utilize FP8 WGMMA appropriately, one must be cognizant of the operand layout restrictions. In the context of a GEMM operation computing $A \times B^{\top}$, where $A$ is an $M\times K$ matrix and $B$ is an $N\times K$ matrix, an operand (either $A$ or $B$) is designated as \emph{mn-major} if its memory arrangement is contiguous along the $M$ or $N$ axis, and as \emph{k-major} if its contiguity is in the $K$ direction. For FP16 WGMMA, operands in SMEM can utilize either mn-major or k-major layout, but for FP8 WGMMA, only k-major organization is permissible. Furthermore, tasks such as attention, which require combining consecutive GEMM operations within a unified kernel, are hindered by mismatched layouts between FP32 accumulators and FP8 operands, thereby complicating the chaining of dependent FP8 WGMMAs.

Within the domain of attention mechanisms, such layout constraints necessitate adjustments to the structure of an FP8 algorithm; the specifics of these modifications are detailed in \cref{sec:algofp8}.

\subsection{Standard Attention and Flash Attention} As described in \citet{dao2022flashattention}, \textbf{standard attention} refers to the typical GPU-based approach to computing attention, where the intermediate matrices $\mathbf{S}$ and $\mathbf{P}$ are explicitly stored in HBM. The essential contribution of \textsc{FlashAttention} is to exploit a localized variant of softmax reduction, thereby sidestepping the costly process of reading from and writing to memory for these intermediates and integrating the attention operation within a unified kernel. This local application of softmax is realized by lines \ref{code-ws:softmax_start}-\ref{code-ws:softmax_end} in the consumer mainloop of \cref{alg:flash3_wgmma_ws_only}, alongside blockwise scaling of $\mathbf{O}$. A clear explanation showing how this scheme faithfully computes $\mathbf{O}$ appears in \cite[\S 2.3.1]{dao2023flashattention2}.

\section{FlashAttention-3: Algorithm}
\label{sec:algo}

This section outlines the \textsc{FlashAttention-3} algorithm. To streamline the exposition, we limit our discussion to the forward pass; details regarding the backward pass can be found in~\cref{sec:algo_ws_bwd}. We begin by detailing the method for combining warp-specialization with a circular SMEM buffer within the foundational structure of \textsc{FlashAttention-2}. Next, we describe how the inherent asynchrony of WGMMA can be utilized to implement a two-stage pipeline that overlaps GEMM and softmax operations. Lastly, we address the adjustments required for FP8 support, discussing both data layout considerations and strategies for preserving accuracy, such as block quantization and incoherent computation.

\subsection{Producer-Consumer asynchrony through warp-specialization and pingpong scheduling}
\label{sec:algo_ws}

\paragraph{Warp-specialization}
Similar to the implementation in \textsc{FlashAttention-2}, the forward pass of \textsc{FlashAttention-3} exhibits substantial parallelism over the dimensions of batch size, head count, and the length of the query sequence. Consequently, it is sufficient to describe the algorithm at the level of a Cooperative Thread Array (CTA), where each CTA processes a particular tile $\mathbf{Q}_i$ from the query matrix to generate the associated output tile $\mathbf{O}_i$. For clarity, we introduce the warp-specialization strategy using a circular SMEM buffer—without considering the additional complexity of overlapping GEMM and softmax computations. Consider the head dimension $d$, the sequence length $N$, and select a query block size $B_r$; this partitions $\mathbf{Q}$ into $T_r = \lceil \frac{N}{B_r} \rceil$ submatrices, denoted as $\mathbf{Q}_1, .., \mathbf{Q}_{T_r}$.

\begin{algorithm}[H]
    \caption{\small\label{alg:flash3_wgmma_ws_only}\textsc{FlashAttention-3} forward pass \textbf{without} intra-consumer overlapping -- CTA view}
    \begin{algorithmic}[1]
\REQUIRE Matrices $\mathbf{Q}_i \in \mathbb{R}^{B_r \times d}$ and $\mathbf{K}, \mathbf{V} \in \mathbb{R}^{N \times d}$ are stored in HBM, with a key block size $B_c$, such that $T_c = \lceil \frac{N}{B_c} \rceil$.
\STATE Instantiate a pipeline control object to handle barrier synchronization with a circular SMEM buffer of $s$ stages.
\IF {executing as producer warpgroup}
\STATE Free up a predetermined number of registers.
\STATE Launch load instruction to bring $\mathbf{Q}_i$ from HBM into shared memory.
\STATE Once loading is done, signal consumer threads that $\mathbf{Q}_i$ is available.
\FOR{$0 \le j < T_c$}
    \STATE Await consumption of the $(j\,\%\,s)$th buffer stage.
    \STATE Execute memory loads for $\mathbf{K}_j$ and $\mathbf{V}_j$ from HBM to shared memory in the $(j\,\%\,s)$th buffer slot.
    \STATE After loading, communicate completion to consumers for $\mathbf{K}_j$ and $\mathbf{V}_j$.
\ENDFOR
\ELSE \STATE Allocate registers in accordance with consumer warp count.
\STATE Locally initialize $\mathbf{O}_i = (0) \in \mathbb{R}^{B_r \times d}$, and set $\ell_i$ and $m_i$ to $(0)$ and $(-\infty)$ respectively, with $\ell_i, m_i \in \mathbb{R}^{B_r}$.
\STATE Wait until $\mathbf{Q}_i$ becomes available in shared memory.
\FOR{$0 \le j < T_c$}
\STATE Block until $\mathbf{K}_j$ is ready in shared memory.
\STATE Calculate $\mathbf{S}_i^{(j)} = \mathbf{Q}_i \mathbf{K}_j^T$ via SS-GEMM, then commit and wait. \label{alg:ws_only_gemm1}
\STATE Save $m_i^{\mathrm{old}} = m_i$ and update $m_i = \mathrm{max}(m_i^{\mathrm{old}}, \mathrm{rowmax}(\mathbf{S}_i^{(j)}))$. \label{code-ws:softmax_start}
\STATE Determine $\widetilde{\mathbf{P}}_i^{(j)} = \mathrm{exp}(\mathbf{S}_i^{(j)} - m_i)$ and update $\ell_i = \mathrm{exp}(m_i^{\mathrm{old}} - m_i)\,\ell_i + \mathrm{rowsum}(\widetilde{\mathbf{P}}_i^{(j)})$. \label{code-ws:softmax_end}
\STATE Block until $\mathbf{V}_j$ arrives in shared memory.
\STATE Update $\mathbf{O}_i = \mathrm{diag}(\mathrm{exp}(m_i^{\mathrm{old}} - m_i))^{-1}\mathbf{O}_i + \widetilde{\mathbf{P}}_i^{(j)}\mathbf{V}_j$ via RS-GEMM; commit and wait. \label{alg:ws_only_gemm2}
\STATE Release the $(j\,\%\,s)$th buffer stage to the producer group.
\ENDFOR
\STATE Finalize computation with $\mathbf{O}_i = \mathrm{diag}(\ell_i)^{-1}\mathbf{O}_i$ and $L_i = m_i + \log(\ell_i)$.
\STATE Output $\mathbf{O}_i$ and $L_i$ to HBM, storing them as the $i$th block of matrices $\mathbf{O}$ and $L$.
\ENDIF
\end{algorithmic}
\end{algorithm}

In our deployment of \cref{alg:flash3_wgmma_ws_only} on the Hopper architecture, resource (de)allocations are managed using \verb|setmaxnreg|. Data loading operations employ TMA, specifically for fetching $\mathbf{Q}_i$ and the collection $\{ \mathbf{K}_j, \mathbf{V}_j \}_{0 \leq j < T_c}$. WGMMA is utilized for carrying out the GEMM computations inside the consumer mainloop. The prefixes SS and RS are applied to denote whether the first operand originates from shared memory or the register file, respectively. To analyze the execution pattern described in \cref{alg:flash3_wgmma_ws_only}, it is important to recognize that TMA load commands are asynchronous—thus, they do not block on previous loads finishing. Furthermore, during the producer mainloop’s first $s$ cycles, there will be no waiting instructions, as this phase is dedicated to filling the buffer.

\paragraph{Pingpong scheduling} The asynchronous execution characteristics of WGMMA and TMA, in combination with warp specialization, facilitate the possibility of interleaving the softmax processing for one warpgroup with the GEMM computation for another. To illustrate the significance of this approach, it is important to recognize that on contemporary hardware accelerators, non-matmul tasks achieve considerably lower throughput than matmul tasks. For instance, the H100 SXM5 GPU delivers 989 TFLOPS for FP16 matmul operations but is limited to only 3.9 TFLOPS for special function operations such as exponentiation\footnote{According to the CUDA programming guide, each streaming multiprocessor (SM) is capable of executing 16 special function operations per clock cycle. By multiplying 16 with 132 SMs and the 1830 MHz clock frequency, the result is 3.9 TFLOPS of special functions.}, which are essential for softmax computation. In the context of the attention forward pass with FP16 precision and a head dimension of 128, the number of matmul FLOPS is 512 times greater than the number of exponential operations, yet the exponential throughput is 256 times lower, resulting in exponentials occupying up to half the runtime relative to matmul. This imbalance becomes even more pronounced with FP8, where the matmul throughput increases twofold, while the exponential throughput remains unchanged.

The exponential operation is handled by a dedicated hardware component (the multi-function unit), so, in an optimal scenario, it is preferable to overlap the exponential computation with the matrix multiplication executed on the Tensor Cores. To achieve this, synchronization barriers (\texttt{bar.sync} instructions) are inserted to enforce an execution order: the GEMMs (GEMM1, which is $\mathbf{P} \mathbf{V}$ for one iteration, and GEMM0, which is $\mathbf{Q} \mathbf{K}^\top$ for the subsequent iteration) from warpgroup 1 are scheduled before those of warpgroup 2. Consequently, while warpgroup 1 is carrying out the softmax computation, warpgroup 2 proceeds with its GEMMs, and then the tasks alternate—warpgroup 2 performs softmax as warpgroup 1 works on GEMMs. This back-and-forth approach, known as ``pingpong'' scheduling, is depicted in~\cref{fig:pingpong_scheduling}. While in practical scenarios this alternation does not occur as seamlessly as illustrated, empirical evidence suggests that it yields notable performance gains (for example, throughput can increase from 570 TFLOPS to 620-640 TFLOPS for FP16 forward passes with a head dimension of 128 and sequence length 8192).

\paragraph{Attention variants} In the case of multi-query attention~\citep{shazeer2019fast} and grouped query attention~\citep{ainslie2023gqa}, our methodology aligns with \textsc{FlashAttention-2}: we modify tensor indexing strategies to ensure $\mathbf{K}$ and $\mathbf{V}$ are not replicated within HBM.

\subsection{Intra-warpgroup overlapping GEMMs and softmax}

Within a single warpgroup, it is possible to interleave certain instructions from the softmax operation with those from the GEMMs. Here, we outline a method for achieving this overlap.

In the attention algorithm, operations within the inner loop (main loop) have sequential dependencies that impede parallelization within a single iteration. For example, (local) softmax (lines \ref{code-ws:softmax_start} to \ref{code-ws:softmax_end}) relies on the output $\mathbf{S}_i^{(j)}$ of the first GEMM, while the second GEMM takes its result $\widetilde{\mathbf{P}}_i^{(j)}$ as an operand. Indeed, the wait statements in lines \ref{alg:ws_only_gemm1} and \ref{alg:ws_only_gemm2} of \cref{alg:flash3_wgmma_ws_only} serialize the execution of softmax and GEMMs. However, we can break these dependencies by pipelining across iterations through additional buffers in registers. Pursuing this idea, we propose the following two-stage\footnote{Note that the number of stages of the overlapping scheme is bounded by, but need not equal, the number $s$ of stages in the circular SMEM buffer.} GEMM-softmax pipelining algorithm:

\begin{algorithm}[H]
    \caption{\small\label{alg:flash3_wgmma}\textsc{FlashAttention-3} consumer warpgroup forward pass}
    \begin{algorithmic}[1]
      \REQUIRE  Matrices $\mathbf{Q}_i \in \mathbb{R}^{B_r \times d}$ and $\mathbf{K}, \mathbf{V} \in \mathbb{R}^{N \times d}$ stored in HBM, key block size $B_c$ with $T_c = \lceil \frac{N}{B_c} \rceil$.
      \STATE Adjust register allocation dynamically depending on the number of consumer warps.
      \STATE Allocate and set $\mathbf{O}_i = (0) \in \mathbb{R}^{B_r \times d}$, as well as $\ell_i, m_i = (0), (-\infty) \in \mathbb{R}^{B_r}$, in on-chip memory.
      \STATE Wait for $\mathbf{Q}_i$ and $\mathbf{K}_0$ to become available in shared memory.
      \STATE Execute WGMMA to evaluate $\mathbf{S}_{\mathrm{cur}} = \mathbf{Q}_i \mathbf{K}_0^T$, commit results, and synchronize.
      \STATE Free the buffer's $0$th stage for $\mathbf{K}$.
      \STATE Compute $m_{i}$, $\tilde{\mathbf{P}}_{\mathrm{cur}}$, and $\ell_{i}$ from $\mathbf{S}_{\mathrm{cur}}$ and apply rescaling to $\mathbf{O}_i$.
      \FOR{$1 \le j < T_c - 1$}
        \STATE Wait for transfer of $\mathbf{K}_j$ into shared memory.
        \label{code:mainloop_start}
        \STATE Invoke WGMMA for $\mathbf{S}_{\mathrm{next}} = \mathbf{Q}_i \mathbf{K}_{j}^T$, commit but do not block.
        \label{code:first_wgmma}
        \STATE Await loading of $\mathbf{V}_{j-1}$ into shared memory.
        \STATE Update $\mathbf{O}_{i} = \mathbf{O}_{i} + \tilde{\mathbf{P}}_{\mathrm{cur}} \mathbf{V}_{j-1}$ via WGMMA, commit and continue without waiting.
        \label{code:second_wgmma}
        \STATE Wait for completion of WGMMA calculation for $\mathbf{Q}_i \mathbf{K}_{j}^T$.
        \STATE Derive $m_{i}$, $\tilde{\mathbf{P}}_{\mathrm{next}}$, and $\ell_{i}$ using $\mathbf{S}_{\mathrm{next}}$. 
        \label{code:softmax}
        \STATE Wait until WGMMA for $\tilde{\mathbf{P}}_{\mathrm{cur}} \mathbf{V}_{j-1}$ is finished, then rescale $\mathbf{O}_i$.
        \STATE Release stage $(j\,\%\,s)$ in the $\mathbf{K}$ buffer, and $(j-1\,\%\,s)$ in the $\mathbf{V}$ buffer.
        \STATE Assign $\mathbf{S}_{\mathrm{cur}} \leftarrow \mathbf{S}_{\mathrm{next}}$.
        \label{code:mainloop_end}
      \ENDFOR
      \STATE Wait until $\mathbf{V}_{T_c - 1}$ is copied to shared memory.
      \STATE Compute 
      $\mathbf{O}_{i} = \mathbf{O}_{i} + \tilde{\mathbf{P}}_{\mathrm{last}} \mathbf{V}_{T_c - 1}$ using WGMMA. Commit and synchronize.

\STATE Epilogue: Adjust $\mathbf{O}_{i}$ using $m_{i}$. Determine $L_{i}$ using both $m_{i}$ and $\ell_{i}$.
Store $\mathbf{O}_{i}$ and $L_{i}$ in HBM, assigning them as the $i$-th block of $\mathbf{O}$ and $L$, respectively.
\end{algorithmic}
\end{algorithm}

\cref{alg:flash3_wgmma} serves as a substitute for the consumer pathway in \cref{alg:flash3_wgmma_ws_only}, thereby constituting the full \textsc{FlashAttention-3} algorithm in FP16 format. Conceptually, we denote WGMMA as shorthand for asynchronous GEMM computations. During the primary loop (spanning lines \ref{code:mainloop_start} through \ref{code:mainloop_end}), the execution of the second WGMMA operation in the current iteration $j$ (at line \ref{code:second_wgmma}) is performed concurrently with the softmax computations initiated in the following iteration $j+1$ (specified at line \ref{code:softmax}).

Although the pipelined design depicted above suggests substantial theoretical advantages in throughput, several real-world limitations must be addressed:

\paragraph{Compiler reordering} Although the pseudocode assumes a certain order of operations, the NVCC compiler tends to optimize instruction sequences by reordering them. This optimization can interfere with the intended pipelining of WGMMA and accompanying instructions, potentially reducing the anticipated performance gains or causing unintended execution artifacts. Examination of SASS output reveals that the compiler does achieve the desired instruction overlap, as detailed in Section~\ref{sec:2-stage-sass}.

\paragraph{Register pressure} Efficient execution relies on minimizing register spills. The dual-stage pipeline inherently increases register usage because temporary values and contextual state for each pipeline stage must be stored, notably the additional $\mathbf{S}_{\mathrm{next}}$ for stage transitions. Consequently, each threadblock incurs extra register consumption amounting to $B_r \times B_c \times \text{sizeof}(\text{float})$. This heightened demand could constrain the utilization of larger threadblocks, which also require significant registers, necessitating careful profiling and tuning to optimize resource allocation.

\paragraph{3-stage pipelining} We introduce a 3-stage variant that expands on the previous 2-stage algorithm to overlap a second WGMMA operation with softmax computation. This model may enable even better exploitation of Tensor Core resources, but it increases register requirements due to the added pipeline stage, complicating the balance between tile dimensions and pipeline depth. For further details on the implementation and benchmarking of the 3-stage algorithm, refer to~\cref{sec:3-stage}.

\subsection{Low-precision with FP8}
\label{sec:algofp8}

\textbf{Efficiency: layout transformations.} Executing the forward pass of \textsc{FlashAttention-3} using FP8 precision introduces unique difficulties related to layout conformity, which are not present when employing FP16.

Initially, it is important to recognize that the input tensors $\mathbf{Q}$, $\mathbf{K}$, and $\mathbf{V}$ are usually arranged so that they are contiguous along the head dimension. However, in order to comply with the k-major constraint required by FP8 WGMMA in the second GEMM operation, it is necessary for $\mathbf{V}$—specifically, the sub-tiles of $\mathbf{V}$ that are transferred to SMEM—to be contiguous along the sequence length axis. Since the TMA load mechanism is unable to alter the dimension of contiguity, this necessitates either (1) transposing $\mathbf{V}$ within GMEM beforehand, or (2) performing an in-kernel transpose on the tiles of $\mathbf{V}$ post-transfer to SMEM. For option (1), the transpose can be integrated either (1a) as a fused operation in the epilogue of a prior step such as rotary embedding, or (1b) by invoking a dedicated transpose kernel during pre-processing\footnote{A high-performance transpose kernel can approach device bandwidth utilization~\citep{colfax_cutlass_transpose_2024}.} to swap the sequence length and head strides. Nevertheless, (1a) poses significant challenges for inclusion in standardized libraries, and (1b) is not efficient in memory-constrained contexts such as inference.

For FP8 \textsc{FlashAttention-3}, we select the second approach. When performing the in-kernel transpose, we leverage the LDSM (\verb|ldmatrix|) and STSM (\verb|stmatrix|) instructions. These enable an entire thread warp to cooperatively move data between shared memory (SMEM) and register memory (RMEM) at 128-byte increments.\footnote{According to the PTX documentation, LDSM/STSM operate on $8 \times 8$ matrices with 16-bit values \cite[\S 9.7.13.4.15-16]{ptx}. We, however, exploit their functionality in the context of FP8 by packing pairs of 8-bit entries. Since the transpose modes of LDSM/STSM are unable to separate packed 8-bit values, further data shuffling between LDSM and STSM at the register level becomes necessary to realize a proper tile-wise transpose; for brevity, the specifics are omitted here.} Because LDSM and STSM use registers efficiently, they can be utilized by the producer warpgroup and support data layout transposition during memory operations. Additionally, starting from the second iteration, we schedule the transpose of the upcoming $\mathbf{V}$ tile so that it can proceed concurrently—masked by the two WGMMAs dependent on the prior $\mathbf{V}$ and present $\mathbf{K}$ tiles.

Second, we note that, in contrast to FP16, the memory representation of the FP32 accumulator produced by FP8 WGMMA diverges from the format used by operand A when stored in registers. Portions of both layouts are shown in \cref{fig:rmem_accum} and \cref{fig:rmem_operand}, illustrating the register ordering per thread. Leveraging byte permutation operations, we can reformat the accumulator resulting from the initial WGMMA so that it aligns both with the expectations of the subsequent WGMMA and with the structure of the $\mathbf{V}$ tile, which is generated through in-kernel transpose. Explicitly, referring to \cref{fig:rmem_accum}, we reorder elements as follows:
$$\{ \verb|d0 d1 d4 d5 d2 d3 d6 d7| \},$$
and this pattern repeats for every 8 bytes. Conceptually, this alters the column arrangement of the $\mathbf{P}$ tile, so columns $0189$ now occupy the first four positions. To ensure WGMMA produces the correct result tile, we adjust the in-kernel transpose so that the $\mathbf{V}$ tile is written with a corresponding row permutation.\footnote{By incorporating the in-kernel transpose in this way, we avoid needing shuffle instructions to redistribute register ownership among threads—a requirement described previously in~\citep{colfax_fp8_flashattention_2024}.}

\textbf{Accuracy: block quantization and incoherent processing.} The FP8 (e4m3) numerical format encodes numbers using just 3 mantissa bits alongside 4 exponent bits, inherently introducing greater rounding and quantization error relative to FP16/BF16 formats. Large-scale models are further challenged by the presence of value outliers~\citep{dettmers2208llm, sun2024massive}, where some elements exhibit much larger magnitudes than the rest, which complicates precise quantization. Standard practice involves per-tensor scaling~\citep{micikevicius2022fp8}, assigning a single scalar value for each tensor (such as $\mathbf{Q}$, $\mathbf{K}$, and $\mathbf{V}$) to adjust dynamic range appropriately. To address and minimize FP8-induced inaccuracies in attention calculations, we incorporate two methods:

\begin{enumerate}

\item \textbf{Block quantization}: this approach involves assigning a single scalar per block, partitioning the tensors $\mathbf{Q}$, $\mathbf{K}$, and $\mathbf{V}$ into segments of dimension $B_r \times d$ or $B_c \times d$, and applying quantization to each block independently. The quantization step can be efficiently combined with preceding operations, such as rotary embedding, without incurring additional latency, given that these operations are primarily constrained by memory bandwidth. Since the \textsc{FlashAttention-3} algorithm inherently processes blocks, we can rescale each block of $\mathbf{S}$ to accommodate block quantization without incurring extra computational overhead.

\item \textbf{Incoherent processing}: to mitigate the impact of extreme values, both $\mathbf{Q}$ and $\mathbf{K}$ are premultiplied by a random orthogonal matrix $\mathbf{M}$ prior to FP8 quantization. Owing to the orthogonality property, $\mathbf{M} \mathbf{M}^\top = I$, which ensures that $(\mathbf{Q} \mathbf{M}) (\mathbf{K} \mathbf{M})^\top = \mathbf{Q} \mathbf{K}^\top$, and thus the final attention calculation remains unaffected. This transformation redistributes the outlier values, since each element in $\mathbf{Q} \mathbf{M}$ or $\mathbf{K} \mathbf{M}$ represents a random weighted sum from the original tensors, thereby diminishing quantization artifacts. Following the methodology of \citet{chee2024quip} and~\citet{tseng2024quip}, $\mathbf{M}$ is constructed as a product of random diagonal matrices (with entries $\pm 1$) and a Hadamard matrix, facilitating multiplication in $O(d \log d)$ time rather than $O(d^2)$, and allowing integration with rotary embedding at no extra computational cost.

We confirm in \cref{sec:numerical_error} that both methods achieve up to a 2.6$\times$ reduction in numerical error.

\section{Empirical Validation}
\label{sec:exp}

To implement \textsc{FlashAttention-3}, we employ foundational constructs from CUTLASS~\citep{Thakkar_CUTLASS_2023}, including the WGMMA and TMA abstractions, which we subsequently assess for both performance and precision.

\begin{itemize}

\item \textbf{Attention benchmarking.} We evaluate the performance of \textsc{FlashAttention-3} over a range of sequence lengths and compare its execution time to a baseline PyTorch version, \textsc{FlashAttention-2}, the Triton implementation of \textsc{FlashAttention-2} (optimized for H100 architectures), as well as the cuDNN-based \textsc{FlashAttention-2} provided by the hardware vendor for H100 GPUs. Our results show that \textsc{FlashAttention-3} delivers up to a 2.0$\times$ speedup over \textsc{FlashAttention-2}, and achieves a 1.5$\times$ improvement relative to the Triton variant. Peak performance reaches 740 TFLOPs/s for \textsc{FlashAttention-3}, accounting for 75\% of the theoretical peak throughput of the H100 GPUs.

\item \textbf{Ablation analysis.} Through systematic removal of our design optimizations, we demonstrate that warp-specialization and GEMM-softmax pipelining are key factors in the acceleration observed with \textsc{FlashAttention-3}.

\item \textbf{FP8 attention accuracy.} Our experiments verify that applying block quantization alongside incoherent computation strategies decreases the numerical errors in the FP8 format of \textsc{FlashAttention-3} by a factor of 2.6$\times$.
\end{itemize}

\subsection{Benchmarking Attention}
\label{subsec:benchmark_attn}

We evaluate the execution time of various attention mechanisms on an H100 80GB SXM5 GPU, testing configurations with and without causal masks, and head dimensions of 64 or 128, using FP16 data types. The outcomes, displayed in~\cref{fig:benchmark_attn_fwd} and~\cref{fig:benchmark_attn_bwd}, indicate that \textsc{FlashAttention-3} achieves speedups of approximately 1.5-2.0$\times$ over \textsc{FlashAttention-2} in the forward computation, and between 1.5-1.75$\times$ in the backward computation. Relative to a conventional attention baseline, \textsc{FlashAttention-3} demonstrates improvements of up to 3-16$\times$. Notably, for sequence lengths of 1k or more, \textsc{FlashAttention-3} even outperforms a proprietary library (cuDNN -- closed source) specifically optimized for H100 GPUs.

\paragraph{Benchmark settings:} We vary the sequence length as 512, 1k, ..., 16k, and set batch size so that the total number of tokens is 16k. We set the hidden dimension to 2048, and head dimension to be either 64, 128, or 256 (i.e., 32 heads, 16 heads, or 8 heads). To calculate the FLOPs of the forward pass, we use:

\begin{equation*}
  4 \cdot \text{seqlen}^2 \cdot \text{head dimension} \cdot \text{number of heads}.
\end{equation*}

When using causal masking, the total is halved to reflect that roughly 50% of the computations are performed. The FLOPs for the backward pass are estimated by multiplying the forward pass FLOPs by 2.5; this is based on the presence of 2 matrix multiplications in the forward pass compared to 5 in the backward pass, as a result of recomputation.

Additionally, we evaluate the forward pass runtime using FP8 in comparable environments. The outcomes for a headdim of 256 are illustrated in ~\cref{fig:benchmark_attn_fp8}, with comprehensive results detailed in \cref{sec:benchmark_attn_fp8_full}.

\subsection{Ablation Study: 2-Stage Pipelining Experiments}
\label{sec:2-stage-pipelining-experiments}

We evaluate the impact of disabling both the two-stage WGMMA-softmax pipelining and warp-specialization in the context of non-causal FP16 \textsc{FlashAttention-3}, while maintaining constant parameters $\{\text{batch}, \text{seqlen}, \text{nheads}, \text{hdim}\} = \{ 4, 8448, 16, 128\}$. As presented in~\cref{table:ablation_pipelining}, our findings demonstrate that integrating algorithmic enhancements—specifically, asynchronous processing with warp-specialization and concurrent execution of GEMM and softmax computations—yields notable performance gains, boosting throughput from 570 to 661 TFLOPs.

\subsection{Numerical Error Validation}
\label{sec:numerical_error}

Given the growing concern over the numerical errors inherent in \textsc{FlashAttention}~\citep{golden2024flash}, we perform a comparative analysis between \textsc{FlashAttention-2}, \textsc{FlashAttention-3}, and a conventional attention mechanism, using an FP64-based reference implementation. In order to emulate extreme feature values and activations found in LLMs~\citep{dettmers2208llm, sun2024massive}, we sample the elements of $\mathbf{Q}, \mathbf{K}, \mathbf{V}$ according to the distribution described below:

\begin{equation*}
  \mathcal{N}(0, 1) + \mathcal{N}(0, 100) \cdot \mathrm{Bernoulli}(0.001).
\end{equation*}

Specifically, each element follows a normal distribution with a mean of zero and a standard deviation of 1; however, for 0.1\% of elements, we introduce an additional independent noise term drawn from a normal distribution with standard deviation 10. The resulting root mean squared error (RMSE) is reported in~\cref{table:numerical_error}. In FP16, both \textsc{FlashAttention-2} and \textsc{FlashAttention-3} attain an RMSE that is 1.7$\times$ lower than the conventional implementation, attributable to the storage of intermediate values (specifically the softmax) in FP32 precision. For the FP8 baseline attention, per-tensor scaling is employed; matrix multiplications accumulate results in FP32, while intermediate softmax outcomes are preserved in FP16. Leveraging block quantization and incoherent processing, \textsc{FlashAttention-3} in FP8 delivers 2.6$\times$ higher accuracy relative to this baseline.

\section{Dicussion, Limitations, Conclusion}
\label{sec:discussion}

Our development of \textsc{FlashAttention-3} illustrates that innovative programming paradigms and advances in hardware—including asynchrony and reduced precision—can significantly enhance both the speed and reliability of attention mechanisms. Relative to \textsc{FlashAttention-2}, our implementation achieves attention operations at a rate 1.5-2.0$\times$ faster, while also diminishing FP8 numerical errors by a factor of 2.6 compared to conventional per-tensor quantization methods. There remain several limitations that we aim to resolve in subsequent work: improving performance for LLM inference, incorporating a persistent kernel approach into the FP8 variant,\footnote{Our benchmarking indicates that FP16 \textsc{FlashAttention-3} leverages a persistent kernel as well as load balancing, in contrast to FP8 \textsc{FlashAttention-3}, which lacks this optimization. This contributes to the observation that FP8 \textsc{FlashAttention-3} underperforms for short sequences and causal masking when compared to FP8 cuDNN kernels.} and further investigating the implications of low-precision attention for training at scale. Although our evaluation centers on Hopper GPUs, we anticipate that these methodologies will generalize to other accelerator architectures. Ultimately, we believe that enhanced, efficient attention primitives could enable new possibilities in domains requiring extended context.

\end{document}